\documentclass{book}
\usepackage{fullpage}
\usepackage{svg}
\usepackage{amsmath}
\usepackage{graphicx}

\usepackage{listings}
\usepackage{xcolor}
\usepackage{float}
\definecolor{terminalgreen}{HTML}{8AE234}

% 1. Define Gnuplot language syntax
\lstdefinelanguage{gnuplot}{
    morekeywords={plot, splot, set, unset, title, xlabel, ylabel, zlabel, 
                  output, terminal, using, with, lines, linespoints, points, 
                  style, data, format, xrange, yrange, trange, urange, vrange, 
                  autoscale, grid, key, border, arrow, label, multiplot, replot,
                  fit, via, load, save, cd, pwd, exit, quit},
    sensitive=true,
    morecomment=[l]{\#},
    morestring=[b]",
    morestring=[b]'
}

\lstdefinelanguage{fitlog}{
    sensitive=true,
    morekeywords={Final, set, of, parameters, Asymptotic, Standard, Error}
}

% 2. Define the visual theme for your code blocks
\lstset{
    language=gnuplot,
    basicstyle=\ttfamily\small,      
    keywordstyle=\color{blue}\bfseries, 
    commentstyle=\color{gray}\itshape,  
    stringstyle=\color{purple},      
    numbers=left,                    
    numberstyle=\tiny\color{gray},   
    stepnumber=1,                    
    tabsize=4,                       
    showstringspaces=false,          
    breaklines=true,                 
    frame=none
}

\lstdefinestyle{fitlog}{
    language=fitlog,
    basicstyle=\ttfamily\footnotesize\color{black},
    keywordstyle=\color{terminalgreen}\bfseries,
    commentstyle=\color{terminalcyan},
    numbers=none,
    showstringspaces=false,
    keepspaces=true,
    breaklines=true,
    columns=fullflexible,
    frame=single,
    framerule=0.4pt
}

\lstdefinestyle{datafile}{
    language={},
    basicstyle=\ttfamily\footnotesize\color{magenta},
    numbers=none,
    showstringspaces=false,
    keepspaces=true,
    breaklines=true,
    columns=fullflexible,
    frame=single,
    framerule=0.4pt
}



\begin{document}
\begin{titlepage}
    \centering
    \vspace*{2cm}
    {\LARGE GNU-Plot-Lab-Notebook\par}
    \vspace{3cm}
    \begin{tabular}{rl}
        Roll Number: & \rule{8cm}{0.4pt}\\[1.5cm]
        Name: & Mrinmoy Dey\\[1.5cm]
        Semester: & \rule{8cm}{0.4pt}
    \end{tabular}
\end{titlepage}
%Question (1:xxi)

\newpage
\section*{Section 1}
\addcontentsline{toc}{section}{Section 1}

\subsection*{Question 1}
\addcontentsline{toc}{subsection}{Question 1}
\noindent\textbf{Problem Set 1 (xvii).} Define the piecewise function
\[
f(x)=\begin{cases}x^2,&x<1,\\x,&1\le x\le2,\\x^3/4,&x>2.\end{cases}
\]
Plot it for $0\le x\le3$ and $0\le y\le8$, with $x$-ticks every $0.5$, $y$-ticks every $1$, and the title \textit{Polynomials of x}.

\noindent\textbf{Sol:} The GNUplot script and its output are shown below.
\lstinputlisting[language=gnuplot]{../../Lab-files/GP-Files/1/Q_xvii.gp}
\begin{figure}[H]
    \centering
    \includesvg[width=\textwidth]{../../Lab-files/Output/1/Q_xvii_ou.svg}
    \caption{Output for Problem Set 1 (xvii).}
\end{figure}
\newpage

\subsection*{Question 2}
\addcontentsline{toc}{subsection}{Question 2}
\noindent\textbf{Problem Set 1 (xv).} Plot $y=x\sin(3x)$ for $-3\le x\le3$ and $-4\le y\le4$. Show both axes, label them \texttt{x values} and \texttt{y values}, and set the $x$- and $y$-tick intervals to $0.5$ and $1$, respectively.

\noindent\textbf{Sol:} The GNUplot script and its output are shown below.
\lstinputlisting[language=gnuplot]{../../Lab-files/GP-Files/1/Q_xv.gp}
\begin{figure}[H]
    \centering
    \includesvg[width=\textwidth]{../../Lab-files/Output/1/Q_xv_ou.svg}
    \caption{Output for Problem Set 1 (xv).}
\end{figure}
\newpage

\subsection*{Question 3}
\addcontentsline{toc}{subsection}{Question 3}
\noindent\textbf{Problem Set 1 (i).} Plot $y=f(x)=x^5+2x^2+6$ for $-10\le x\le10$. Show both axes, label them \texttt{x values} and \texttt{new function}, and set the $x$-tick interval to $2$.

\noindent\textbf{Sol:} The GNUplot script and its output are shown below.
\lstinputlisting[language=gnuplot]{../../Lab-files/GP-Files/1/Q_i.gp}
\begin{figure}[H]
    \centering
    \includesvg[width=\textwidth]{../../Lab-files/Output/1/Q_i_ou.svg}
    \caption{Output for Problem Set 1 (i).}
\end{figure}
\newpage

\subsection*{Question 4}
\addcontentsline{toc}{subsection}{Question 4}
\noindent\textbf{Problem Set 1 (xi).} Plot $y=\ln(x^2)$ for $-2\le x\le2$ and $-3\le y\le3$. Show both axes, label them \texttt{x values} and \texttt{y values}, and set the $x$- and $y$-tick intervals to $0.5$ and $1$, respectively.

\noindent\textbf{Sol:} The GNUplot script and its output are shown below.
\lstinputlisting[language=gnuplot]{../../Lab-files/GP-Files/1/Q_xi.gp}
\begin{figure}[H]
    \centering
    \includesvg[width=\textwidth]{../../Lab-files/Output/1/Q_xi_ou.svg}
    \caption{Output for Problem Set 1 (xi).}
\end{figure}
\newpage

\subsection*{Question 5}
\addcontentsline{toc}{subsection}{Question 5}
\noindent\textbf{Problem Set 1 (vi).} Plot $f(x)=2x^3+9x-11$ for $-15\le x\le10$ and $-7000\le y\le2200$. Show the $y$-axis, label the axes \texttt{x value} and \texttt{new function}, and set the $x$-tick interval to $2.5$.

\noindent\textbf{Sol:} The GNUplot script and its output are shown below.
\lstinputlisting[language=gnuplot]{../../Lab-files/GP-Files/1/Q_vi.gp}
\begin{figure}[H]
    \centering
    \includesvg[width=\textwidth]{../../Lab-files/Output/1/Q_vi_ou.svg}
    \caption{Output for Problem Set 1 (vi).}
\end{figure}
\newpage

\subsection*{Question 6}
\addcontentsline{toc}{subsection}{Question 6}
\noindent\textbf{Problem Set 1 (v).} Plot the piecewise functions
\[
f_1(x)=x\;(x>4),\qquad f_2(x)=-x^2\;(-4<x<4),\qquad f_3(x)=-x\;(x<-4)
\]
for $-10\le x\le10$ and $-4\le y\le10$. Set the $x$-tick interval to $1$ and label the axes \texttt{x values} and \texttt{piecewise continuous functions}.

\noindent\textbf{Sol:} The GNUplot script and its output are shown below.
\lstinputlisting[language=gnuplot]{../../Lab-files/GP-Files/1/Q_v.gp}
\begin{figure}[H]
    \centering
    \includesvg[width=\textwidth]{../../Lab-files/Output/1/Q_v_ou.svg}
    \caption{Output for Problem Set 1 (v).}
\end{figure}
\newpage

\subsection*{Question 7}
\addcontentsline{toc}{subsection}{Question 7}
\noindent\textbf{Problem Set 1 (iv).} Plot the piecewise function
\[
f(x)=\begin{cases}x,&-2\le x\le2,\\0.5,&\text{otherwise}\end{cases}
\]
for $-10\le x\le10$ and $-5\le y\le5$. Set the $x$- and $y$-tick intervals to $2$ and $1$, respectively, and label the axes \texttt{x value} and \texttt{jerk function}.

\noindent\textbf{Sol:} The GNUplot script and its output are shown below.
\lstinputlisting[language=gnuplot]{../../Lab-files/GP-Files/1/Q_iv.gp}
\begin{figure}[H]
    \centering
    \includesvg[width=\textwidth]{../../Lab-files/Output/1/Q_iv_ou.svg}
    \caption{Output for Problem Set 1 (iv).}
\end{figure}
\newpage
\section*{Section 2}
\addcontentsline{toc}{section}{Section 2}

\subsection*{Question 1}
\addcontentsline{toc}{subsection}{Question 1}
\noindent\textbf{Problem Set 2 (ii).} Fit the data with $y=mx+c$ over $-7\le x\le7$. Plot the points and fitted line, show the $y$-axis, and choose the $y$-range so $c$ can be read from that axis. Report $m$ and $c$.
\begin{table}[H]
\centering
\caption{Data for Problem Set 2 (ii).}
\begin{tabular}{|r|r|}
\hline
$x$ & $y$ \\ \hline
$-5.1$ & $21.8$ \\ \hline
$-3.4$ & $16.12$ \\ \hline
$-0.5$ & $7.05$ \\ \hline
$1.57$ & $0.42$ \\ \hline
$6.2$ & $-14.2$ \\ \hline
\end{tabular}
\end{table}
\noindent\textbf{Data file contents:}
\lstinputlisting[style=datafile]{../../Lab-files/Data/2/Q_ii.txt}
\noindent\textbf{Sol:} The fitted plot and the parameter estimates with their standard errors are shown below.
\lstinputlisting[language=gnuplot]{../../Lab-files/GP-Files/2/GNU-Plot-files/Q_2_ii.gp}
\begin{figure}[H]
    \centering
    \includesvg[width=\textwidth]{../../Lab-files/Output/2/Q_2_ii_ou.svg}
    \caption{Linear fit for Problem Set 2 (ii).}
\end{figure}
\noindent\textbf{Fit log (parameter estimates and standard errors):}
\lstinputlisting[style=fitlog,firstline=29,lastline=32]{../../Lab-files/GP-Files/2/Fit-log/fit_ii.log}
\newpage

\subsection*{Question 2}
\addcontentsline{toc}{subsection}{Question 2}
\noindent\textbf{Problem Set 2 (iii).} Fit the data with $y=mx+c$ over $0\le x\le9$. Plot the points and fitted line, show the $y$-axis, and choose the $y$-range so $c$ can be read from that axis. Report $m$ and $c$.
\begin{table}[H]
\centering
\caption{Data for Problem Set 2 (iii).}
\begin{tabular}{|r|r|}
\hline
$x$ & $y$ \\ \hline
$2.1$ & $-2.5$ \\ \hline
$4.6$ & $7.7$ \\ \hline
$5.7$ & $12.17$ \\ \hline
$6.6$ & $15.9$ \\ \hline
$8.2$ & $22.4$ \\ \hline
\end{tabular}
\end{table}
\noindent\textbf{Data file contents:}
\lstinputlisting[style=datafile]{../../Lab-files/Data/2/Q_iii.txt}
\noindent\textbf{Sol:} The fitted plot and the parameter estimates with their standard errors are shown below.
\lstinputlisting[language=gnuplot]{../../Lab-files/GP-Files/2/GNU-Plot-files/Q_2_iii.gp}
\begin{figure}[H]
    \centering
    \includesvg[width=\textwidth]{../../Lab-files/Output/2/Q_2_iii_ou.svg}
    \caption{Linear fit for Problem Set 2 (iii).}
\end{figure}
\noindent\textbf{Fit log (parameter estimates and standard errors):}
\lstinputlisting[style=fitlog,firstline=29,lastline=32]{../../Lab-files/GP-Files/2/Fit-log/fit_iii.log}
\newpage

\subsection*{Question 3}
\addcontentsline{toc}{subsection}{Question 3}
\noindent\textbf{Problem Set 2 (iv).} Fit the data with $f(x)=mx+c$ over $-3\le x\le4$. Plot the points and fitted line, show both axes, and choose the $y$-range so $c$ can be read from the $y$-axis. Report $m$ and $c$.
\begin{table}[H]
\centering
\caption{Data for Problem Set 2 (iv).}
\begin{tabular}{|r|r|}
\hline
$x$ & $y$ \\ \hline
$-2.0$ & $-4.15$ \\ \hline
$-1.3$ & $-2.40$ \\ \hline
$1.8$ & $6.60$ \\ \hline
$2.5$ & $8.98$ \\ \hline
$3.6$ & $11.94$ \\ \hline
\end{tabular}
\end{table}
\noindent\textbf{Data file contents:}
\lstinputlisting[style=datafile]{../../Lab-files/Data/2/Q_iv.txt}
\noindent\textbf{Sol:} The fitted plot and the parameter estimates with their standard errors are shown below.
\lstinputlisting[language=gnuplot]{../../Lab-files/GP-Files/2/GNU-Plot-files/Q_2_iv.gp}
\begin{figure}[H]
    \centering
    \includesvg[width=\textwidth]{../../Lab-files/Output/2/Q_2_iv_ou.svg}
    \caption{Linear fit for Problem Set 2 (iv).}
\end{figure}
\noindent\textbf{Fit log (parameter estimates and standard errors):}
\lstinputlisting[style=fitlog,firstline=29,lastline=32]{../../Lab-files/GP-Files/2/Fit-log/fit_iv.log}
\newpage

\subsection*{Question 4}
\addcontentsline{toc}{subsection}{Question 4}
\noindent\textbf{Problem Set 2 (v).} Fit the data with $f(x)=ax^b$ over $0\le x\le3$ and $0\le y\le35$. Plot the points and fitted curve, use the title \textit{Power series fit}, and report $a$ and $b$.
\begin{table}[H]
\centering
\caption{Data for Problem Set 2 (v).}
\begin{tabular}{|r|r|}
\hline
$x$ & $y$ \\ \hline
$0.5$ & $0.21$ \\ \hline
$0.83$ & $1.05$ \\ \hline
$1.16$ & $2.40$ \\ \hline
$1.82$ & $9.01$ \\ \hline
$2.81$ & $30.75$ \\ \hline
\end{tabular}
\end{table}
\noindent\textbf{Data file contents:}
\lstinputlisting[style=datafile]{../../Lab-files/Data/2/Q_v.txt}
\noindent\textbf{Sol:} The fitted plot and the parameter estimates with their standard errors are shown below.
\lstinputlisting[language=gnuplot]{../../Lab-files/GP-Files/2/GNU-Plot-files/Q_2_v.gp}
\begin{figure}[H]
    \centering
    \includesvg[width=\textwidth]{../../Lab-files/Output/2/Q_2_v_ou.svg}
    \caption{Power-series fit for Problem Set 2 (v).}
\end{figure}
\noindent\textbf{Fit log (parameter estimates and standard errors):}
\lstinputlisting[style=fitlog,firstline=29,lastline=32]{../../Lab-files/GP-Files/2/Fit-log/fit_v.log}
\newpage

\subsection*{Question 5}
\addcontentsline{toc}{subsection}{Question 5}
\noindent\textbf{Problem Set 2 (vi).} Fit the data with $f(x)=mx+c$ over $-3\le x\le3$. Plot the points and fitted line, show both axes, and choose the $y$-range so $c$ can be read from the $y$-axis. Report $m$ and $c$.
\begin{table}[H]
\centering
\caption{Data for Problem Set 2 (vi).}
\begin{tabular}{|r|r|}
\hline
$x$ & $y$ \\ \hline
$-2.9$ & $-5.96$ \\ \hline
$-1.4$ & $-3.81$ \\ \hline
$0.8$ & $-1.10$ \\ \hline
$1.7$ & $0.21$ \\ \hline
$2.5$ & $1.30$ \\ \hline
\end{tabular}
\end{table}
\noindent\textbf{Data file contents:}
\lstinputlisting[style=datafile]{../../Lab-files/Data/2/Q_vi.txt}
\noindent\textbf{Sol:} The fitted plot and the parameter estimates with their standard errors are shown below.
\lstinputlisting[language=gnuplot]{../../Lab-files/GP-Files/2/GNU-Plot-files/Q_2_vi.gp}
\begin{figure}[H]
    \centering
    \includesvg[width=\textwidth]{../../Lab-files/Output/2/Q_2_vi_ou.svg}
    \caption{Linear fit for Problem Set 2 (vi).}
\end{figure}
\noindent\textbf{Fit log (parameter estimates and standard errors):}
\lstinputlisting[style=fitlog,firstline=29,lastline=32]{../../Lab-files/GP-Files/2/Fit-log/fit_vi.log}
\newpage

\subsection*{Question 6}
\addcontentsline{toc}{subsection}{Question 6}
\noindent\textbf{Problem Set 2 (vii).} Fit the data with $f(x)=ax^2+b$ over $-3\le x\le3$ and $0\le y\le16$. Plot the points and fitted curve, use the title \textit{Polynomial fit}, and report $a$ and $b$.
\begin{table}[H]
\centering
\caption{Data for Problem Set 2 (vii).}
\begin{tabular}{|r|r|}
\hline
$x$ & $y$ \\ \hline
$-2.8$ & $14.65$ \\ \hline
$-1.2$ & $5.81$ \\ \hline
$0.7$ & $4.48$ \\ \hline
$1.8$ & $8.33$ \\ \hline
$2.9$ & $15.56$ \\ \hline
\end{tabular}
\end{table}
\noindent\textbf{Data file contents:}
\lstinputlisting[style=datafile]{../../Lab-files/Data/2/Q_vii.txt}
\noindent\textbf{Sol:} The fitted plot and the parameter estimates with their standard errors are shown below.
\lstinputlisting[language=gnuplot]{../../Lab-files/GP-Files/2/GNU-Plot-files/Q_2_vii.gp}
\begin{figure}[H]
    \centering
    \includesvg[width=\textwidth]{../../Lab-files/Output/2/Q_2_vii_ou.svg}
    \caption{Polynomial fit for Problem Set 2 (vii).}
\end{figure}
\noindent\textbf{Fit log (parameter estimates and standard errors):}
\lstinputlisting[style=fitlog,firstline=29,lastline=32]{../../Lab-files/GP-Files/2/Fit-log/fit_vii.log}
\newpage

\subsection*{Question 7}
\addcontentsline{toc}{subsection}{Question 7}
\noindent\textbf{Problem Set 2 (viii).} Fit the data with $f(x)=ae^{bx}$. Plot the points and fitted curve, use the title \textit{Data fitting}, label the axes $x$ and $f(x)$, and report $a$ and $b$.
\begin{table}[H]
\centering
\caption{Data for Problem Set 2 (viii).}
\begin{tabular}{|r|r|}
\hline
$x$ & $f(x)$ \\ \hline
$-3.0$ & $3.422$ \\ \hline
$-2.0$ & $4.057$ \\ \hline
$-1.0$ & $4.808$ \\ \hline
$0.0$ & $5.7$ \\ \hline
$1.0$ & $6.756$ \\ \hline
$2.0$ & $8.008$ \\ \hline
$3.0$ & $9.492$ \\ \hline
$4.0$ & $11.251$ \\ \hline
\end{tabular}
\end{table}
\noindent\textbf{Data file contents:}
\lstinputlisting[style=datafile]{../../Lab-files/Data/2/Q_viii.txt}
\noindent\textbf{Sol:} The fitted plot and the parameter estimates with their standard errors are shown below.
\lstinputlisting[language=gnuplot]{../../Lab-files/GP-Files/2/GNU-Plot-files/Q_2_viii.gp}
\begin{figure}[H]
    \centering
    \includesvg[width=\textwidth]{../../Lab-files/Output/2/Q_2_viii_ou.svg}
    \caption{Exponential fit for Problem Set 2 (viii).}
\end{figure}
\noindent\textbf{Fit log (parameter estimates and standard errors):}
\lstinputlisting[style=fitlog,firstline=28,lastline=31]{../../Lab-files/GP-Files/2/Fit-log/fit_viii.log}
\newpage
\section*{Section 3}
\addcontentsline{toc}{section}{Section 3}

\subsection*{Question 1}
\addcontentsline{toc}{subsection}{Question 1}
\noindent\textbf{Problem Set 3 (i).} Show the polar plots of $\sin^3(t)$ and $\cos^2(2t)$. Label the axes ``t value'' and ``the range'', place ``polar plot'' at $(0,-0.7)$, and show the vertical axis.
\noindent\textbf{Sol:} The corresponding Gnuplot program and output are below.
\lstinputlisting[language=gnuplot]{../../Lab-files/GP-Files/3/Q_3_i.gp}
\begin{figure}[htbp]
    \centering
    \includesvg[width=\textwidth]{../../Lab-files/Output/3/Q_3_i_ou.svg}
    \caption{Polar plots of $\sin^3(t)$ and $\cos^2(2t)$.}
\end{figure}

\newpage
\subsection*{Question 2}
\addcontentsline{toc}{subsection}{Question 2}
\noindent\textbf{Problem Set 3 (ii).} Show the polar plots of $\tan(6t)$ and $\cot(6t)$ over ranges $-20$ to $20$ on both axes. Place the label ``tan(6t), cot(6t)'' at $(4,12)$ and show both axes.
\noindent\textbf{Sol:} The corresponding Gnuplot program and output are below.
\lstinputlisting[language=gnuplot]{../../Lab-files/GP-Files/3/Q_3_ii.gp}
\begin{figure}[htbp]
    \centering
    \includesvg[width=\textwidth]{../../Lab-files/Output/3/Q_3_ii_ou.svg}
    \caption{Polar plots of $\tan(6t)$ and $\cot(6t)$.}
\end{figure}

\newpage
\subsection*{Question 3}
\addcontentsline{toc}{subsection}{Question 3}
\noindent\textbf{Problem Set 3 (iii).} Plot the parametric curve $x=2\cos^3t$, $y=2\sin^3t$ for $-\pi\leq t\leq\pi$. Label the axes as specified and title the plot ``Astroid''.
\noindent\textbf{Sol:} The corresponding Gnuplot program and output are below.
\lstinputlisting[language=gnuplot]{../../Lab-files/GP-Files/3/Q_3_iii.gp}
\begin{figure}[htbp]
    \centering
    \includesvg[width=\textwidth]{../../Lab-files/Output/3/Q_3_iii_ou.svg}
    \caption{Astroid parametric plot.}
\end{figure}

\newpage
\subsection*{Question 4}
\addcontentsline{toc}{subsection}{Question 4}
\noindent\textbf{Problem Set 3 (iv).} Plot the polar spiral $r=t$ for $0\leq t\leq6\pi$ using $10000$ samples. Show both axes, label them as specified, and title the plot ``Spiral''.
\noindent\textbf{Sol:} The corresponding Gnuplot program and output are below.
\lstinputlisting[language=gnuplot]{../../Lab-files/GP-Files/3/Q_3_iv.gp}
\begin{figure}[htbp]
    \centering
    \includesvg[width=\textwidth]{../../Lab-files/Output/3/Q_3_iv_ou.svg}
    \caption{Polar spiral $r=t$.}
\end{figure}

\newpage
\subsection*{Question 5}
\addcontentsline{toc}{subsection}{Question 5}
\noindent\textbf{Problem Set 3 (v).} Plot the parametric curve $x=2\sin t$, $y=2\sin(2t)$ for $-\pi\leq t\leq\pi$, with the specified axis labels and title ``Lissajous Figure''.
\noindent\textbf{Sol:} The corresponding Gnuplot program and output are below.
\lstinputlisting[language=gnuplot]{../../Lab-files/GP-Files/3/Q_3_v.gp}
\begin{figure}[htbp]
    \centering
    \includesvg[width=\textwidth]{../../Lab-files/Output/3/Q_3_v_ou.svg}
    \caption{Lissajous parametric plot.}
\end{figure}

\newpage
\subsection*{Question 6}
\addcontentsline{toc}{subsection}{Question 6}
\noindent\textbf{Problem Set 3 (vi).} Plot the polar curves $r_1=4+3\cos t$ and $r_2=3+4\cos t$. Label the axes as specified, place ``polar plot'' at $(-6,6)$, and title the plot ``A pair of Limacon''.
\noindent\textbf{Sol:} The corresponding Gnuplot program and output are below.
\lstinputlisting[language=gnuplot]{../../Lab-files/GP-Files/3/Q_3_vi.gp}
\begin{figure}[htbp]
    \centering
    \includesvg[width=\textwidth]{../../Lab-files/Output/3/Q_3_vi_ou.svg}
    \caption{A pair of limaçons.}
\end{figure}

\end{document}